\section{Denominator algebra versus physical BPS algebra}

The Borcherds correction algebra constructed here is a denominator algebra.
The construction proves a denominator identity and a signed-index equality.
A bracket-level physical BPS algebra is a further object.

\begin{definition}[The four levels of structure]
Let
$$
\alpha:\Gamma_{\mathrm{BPS}}\longrightarrow \La^{2,1}_{II},\qquad
\alpha(n,l,m)=2nf_2-lf_3+2mf_{-2},
$$
so that
$$
q^nr^ls^m=\exp(-\pi i(\alpha(n,l,m),z)).
$$
For $\gamma=(n,l,m)\in\Gamma_{\mathrm{eff}}$ put
$$
I_{\mathrm{BPS}}(\gamma):=f(nm,l).
$$
If $\beta\notin\mathcal R_+$, use the convention
$$
\mathfrak g_\beta=0,\qquad \smult(\beta)=0.
$$
There are four different assertions.

\begin{enumerate}
\item \textbf{Proved: denominator identity.}
The generalized Borcherds--Kac--Moody superalgebra $\autcor$ constructed
from the real simple roots $\delta_1,\delta_2,\delta_3$ and the imaginary
simple-root multiplicities has denominator
$$
\operatorname{den}(\autcor)(z)
=e^{-2\pi i(\rho,z)}
\prod_{\beta\in\mathcal R_+}
\left(1-e^{-2\pi i(\beta,z)}\right)^{\smult(\beta)}
=\frac1{64}\Delta_5(2Z).
$$
This is the Weyl--Kac--Borcherds denominator identity, in the normalization
used here; it is the automorphic correction statement of
Borcherds and Gritsenko--Nikulin, written in Kac denominator form
\cite{KAC:1,B,GN}.

\item \textbf{Proved: signed root superdimensions are BPS-index
coefficients.}
For every $\gamma=(n,l,m)\in\Gamma_{\mathrm{eff}}$,
$$
\smult(\alpha(\gamma))
=\rootmult_{\overline0}(\alpha(\gamma))
-\rootmult_{\overline1}(\alpha(\gamma))
=f(nm,l)=I_{\mathrm{BPS}}(\gamma).
$$
This is an equality of signed dimensions, or Euler characteristics.  It is
not an isomorphism of state spaces.

\item \textbf{Interpretive: Harvey--Moore charge grading.}
Following Harvey--Moore \cite{HM}, the charge lattice is the natural grading
lattice for a candidate algebra of BPS states:
$$
\gamma+\gamma'\quad\hbox{is the degree of a bracket of degrees }
\gamma,\gamma'.
$$
In the present paper this physical principle is used only through the
degree map $\alpha$ and the compatibility
$$
\alpha(\gamma+\gamma')=\alpha(\gamma)+\alpha(\gamma').
$$
Thus $\autcor$ is a charge-graded denominator algebra.  The word
``charge-graded'' records the Harvey--Moore interpretation; it does not
construct the physical bracket on BPS states.

\item \textbf{Not proved here: bracket-level physical BPS algebra.}
A physical BPS algebra would require actual protected state spaces
$$
\mathcal H_{\gamma,\overline0}^{\mathrm{BPS}},\qquad
\mathcal H_{\gamma,\overline1}^{\mathrm{BPS}},
\qquad \gamma\in\Gamma_{\mathrm{eff}},
$$
with
$$
\mathcal H_{\gamma}^{\mathrm{BPS}}
=\mathcal H_{\gamma,\overline0}^{\mathrm{BPS}}
\oplus
\mathcal H_{\gamma,\overline1}^{\mathrm{BPS}},
$$
and bilinear brackets
$$
[-,-]_{\mathrm{BPS}}:
\mathcal H_{\gamma}^{\mathrm{BPS}}\otimes
\mathcal H_{\gamma'}^{\mathrm{BPS}}
\longrightarrow
\mathcal H_{\gamma+\gamma'}^{\mathrm{BPS}}
$$
satisfying the super Jacobi identity and compatible with the physical
charge addition.  It would also require a comparison, at least after
passing through $\alpha$, with the root-space bracket of $\autcor$.
\end{enumerate}
\end{definition}

\begin{proposition}[What the denominator determines]
The identity
$$
\frac1{64}\Delta_5(2Z)
=e^{-2\pi i(\rho,z)}
\prod_{\beta\in\mathcal R_+}
\left(1-e^{-2\pi i(\beta,z)}\right)^{\smult(\beta)}
$$
determines the signed root superdimensions $\smult(\beta)$ appearing in
the product expansion.  It does not determine the separate dimensions
$\rootmult_{\overline0}(\beta)$ and $\rootmult_{\overline1}(\beta)$, and it
does not determine the structure constants
$$
[\mathfrak g_\beta,\mathfrak g_{\beta'}]\longrightarrow
\mathfrak g_{\beta+\beta'}.
$$
\end{proposition}

\begin{proof}
The logarithm of the product determines only the exponent attached to each
monomial $e^{-2\pi i(\beta,z)}$.  That exponent is the integer
$\smult(\beta)$.  For $c\in\Z$, the equality $c=a-b$ has infinitely many
solutions $(a,b)\in\Z_{\geq0}^2$: use $(c+k,k)$ if $c\geq0$ and
$(k,k-c)$ if $c<0$.  Thus the product does not recover the two parity
dimensions separately.  No bracket appears in the product:
the denominator contains the graded Euler characteristics of root spaces,
not the bilinear maps between them.  Hence the denominator is insufficient
to construct a physical BPS-state Lie superalgebra.
\end{proof}

\begin{proposition}[Index equality is not a state-space identification]
The equality
$$
\smult(\alpha(n,l,m))=f(nm,l)
$$
proves that the root-space Euler characteristic of $\autcor$ in degree
$\alpha(n,l,m)$ equals the one-particle BPS index coefficient in charge
$(n,l,m)$.  It does not prove an isomorphism
$$
\mathfrak g_{\alpha(n,l,m)}
\cong
\mathcal H_{(n,l,m)}^{\mathrm{BPS}}
$$
and it does not prove equality of brackets.
\end{proposition}

\begin{proof}
The left-hand side is defined inside the algebraic Borcherds correction
algebra:
$$
\smult(\alpha)=
\dim\mathfrak g_{\alpha,\overline0}
-\dim\mathfrak g_{\alpha,\overline1}.
$$
The right-hand side is the coefficient of the weak Jacobi form
$$
\phi_{0,1}(\tau,z)=\sum_{n\geq0,\;l\in\Z}f(n,l)q^nr^l,
$$
transported to three charges by $\gamma=(n,l,m)\mapsto f(nm,l)$.  Equality
of these integers is exactly what the Borcherds product proves.  An
isomorphism of state spaces would be a linear map preserving the two
gradings and parities.  An isomorphism of algebras would additionally
preserve the brackets.  Neither datum is present in the product formula.
\end{proof}

\begin{definition}[Bracket-level realization]
A bracket-level physical realization of the denominator algebra is the
following extra data.
\begin{enumerate}
\item A $\Gamma_{\mathrm{eff}}$-graded super vector space
$$
\mathcal H_{\mathrm{BPS}}
=\bigoplus_{\gamma\in\Gamma_{\mathrm{eff}}}
\mathcal H_{\gamma}^{\mathrm{BPS}},
\qquad
\mathcal H_{\gamma}^{\mathrm{BPS}}
=\mathcal H_{\gamma,\overline0}^{\mathrm{BPS}}
\oplus
\mathcal H_{\gamma,\overline1}^{\mathrm{BPS}},
$$
with finite-dimensional homogeneous pieces and protected index
$$
\dim\mathcal H_{\gamma,\overline0}^{\mathrm{BPS}}
-\dim\mathcal H_{\gamma,\overline1}^{\mathrm{BPS}}
=f(nm,l)
\quad(\gamma=(n,l,m)).
$$

\item A Lie superbracket on $\mathcal H_{\mathrm{BPS}}$ satisfying
$$
[\mathcal H_\gamma^{\mathrm{BPS}},
\mathcal H_{\gamma'}^{\mathrm{BPS}}]_{\mathrm{BPS}}
\subset
\mathcal H_{\gamma+\gamma'}^{\mathrm{BPS}}.
$$

\item A bracket-preserving comparison with $\autcor$, for example a
graded Lie-superalgebra homomorphism whose degree map is $\alpha$, and
which identifies the denominator root spaces after taking signed
dimensions.
\end{enumerate}
\end{definition}

\begin{remark}[Status]
The present manuscript constructs $\autcor$ as a generalized
Borcherds--Kac--Moody superalgebra and proves its denominator identity.
It also proves that the signed root superdimensions in the positive Weyl
chamber are the Fourier coefficients $f(nm,l)$ of the K3 one-particle
index.  It does not construct the protected state spaces
$\mathcal H_{\gamma}^{\mathrm{BPS}}$ with a physical bracket.  Therefore
the accurate statement is
$$
\autcor\ \hbox{is the Borcherds denominator algebra whose signed root
superdimensions reproduce the BPS index,}
$$
not
$$
\autcor=\hbox{the physical algebra of BPS states}.
$$
\end{remark}

\begin{remark}[Recommended notation]
Use
$$
\mathcal A_{\mathrm{den}}:=\autcor
$$
for the object constructed in the paper.  Reserve
$$
\mathcal A_{\mathrm{BPS}}^{\mathrm{phys}}
$$
for a bracket-level Harvey--Moore BPS algebra if such an algebra is
constructed.  The only unconditional comparison established here is
$$
\operatorname{den}(\mathcal A_{\mathrm{den}})
=\frac1{64}\Delta_5(2Z),
\qquad
\smult_{\mathcal A_{\mathrm{den}}}(\alpha(n,l,m))=f(nm,l).
$$
\end{remark}
